\documentclass[12pt]{article}

\title{CR3BP Overview}
\author{Yannis ABDELOUAHAB}
\date{2025-2026}
\usepackage[a4paper, total={7in,9in}]{geometry}
\usepackage{graphicx}
\usepackage[french]{babel}




\begin{document}
	\maketitle
	\tableofcontents
	
	\begin{abstract}
		Le problème à trois corps restreint (CR3BP en anglais) est l'étude du mouvement relatif de trois masses, dont une de masse négligeable devant les trois autres. Nous chercherons à déterminer les coordonées des Points Lagrange, qui sont cinq points d'équilibre du CR3BP. Nous mènerons ensuite une étude de leur stabilité. A terme, nous chercherons à déterminer des orbites périodiques autour des points L1 et L2 en particulier, ainsi que mener une recherche de trajectoires de transfert qui ferait utilisation des variétés stables et instables liées aux Points Lagrange.
	\end{abstract}
	
	\section{Coordonées des Points Lagrange}
	\subsection{Définition du système}
	Nous étudions le système \emph{satellite de masse m} que l'on associe au point matériel $M$. On se place dans le référentiel galiléen -- que l'on note $R_G$ -- qui à pour centre le barycentre des deux masses $M_1$ et $M_2$ qu'on confondera avec le point matériel auquel leur astre respectif est associé.
	
	On cherche à exprimer les coordonées de $M$ dans le référentiel \textit{non-galiléen} -- $R_{NG}$ -- en rotation dans le référentiel galiléen $R_G$ pour simplifier les équations du mouvement. $R_{NG}$ est alors défini comme le référentiel dans lequel $M_1$ et $M_2$ sont immobiles. On mènera la majorité de notre étude dans ce référentiel.
	
	\subsection{Coordonées dans $R_{NG}$}
	On se place comme 
	
	
\end{document}